\documentclass{article}

\usepackage[T2A]{fontenc}
\usepackage[utf8]{inputenc}
\usepackage{helvet}
\usepackage{geometry}
\usepackage{xcolor}
\usepackage{eso-pic}
\usepackage{fancyhdr}
\usepackage{parskip}
\usepackage{microtype}
\usepackage[english, ukrainian]{babel}
\usepackage{url}
\usepackage{amsmath}
\usepackage{amssymb}
\usepackage{amsthm}
\usepackage{longtable}
\usepackage{booktabs}
\usepackage{hyperref}
\usepackage{titlesec}

\definecolor{nato}{HTML}{293757}
\definecolor{footergray}{gray}{0.65}

\geometry{a4paper,left=3cm,right=3cm,top=3cm,bottom=3cm}
\renewcommand{\familydefault}{\sfdefault}
\setcounter{secnumdepth}{3}

\titleformat{\section}
  {\color{nato}\Huge\bfseries}{\thesection.}{1em}{}
\titlespacing*{\section}{0pt}{30pt}{12pt}

\titleformat{\subsection}
  {\color{nato}\LARGE\bfseries}{\thesubsection.}{1em}{}
\titlespacing*{\subsection}{0pt}{20pt}{8pt}

\titleformat{\subsubsection}
  {\color{nato}\Large\bfseries}{\thesubsubsection.}{1em}{}
\titlespacing*{\subsubsection}{0pt}{10pt}{6pt}

\fancyhf{}
\fancyhead[R]{\small\color{footergray} 25 червня 2026}
\fancyfoot[L]{\small\color{footergray} Регламент проєктного планування та управління ДП «ІСС»}
\fancyfoot[R]{\small\color{footergray}\thepage}

\newcommand\HUGE[1]{\fontsize{40}{40}\selectfont #1}

\renewcommand{\headrulewidth}{0pt}
\renewcommand{\footrulewidth}{0.4pt}
\renewcommand{\footrule}{\color{footergray}\hrule width \textwidth height 0.4pt}

\begin{document}

\pagestyle{fancy}

\begin{titlepage}
  \AddToShipoutPictureBG*{\AtPageLowerLeft{\color{nato}\rule{\paperwidth}{\paperheight}}}
  \color{white}\thispagestyle{empty}\vspace*{3cm}\centering
  {\Huge \textbf{РЕГЛАМЕНТ ПРОЄКТНОГО ПЛАНУВАННЯ} \par} \vspace{1cm}
  {\HUGE \textbf{ДП «Інформаційні судові системи»} \par} \vspace{2cm}
  {\LARGE \textbf{Узгоджений регламент розробки програмного забезпечення \\ та супроводу інформаційних систем} \par} \vspace{1.5cm}
  {\large\textbf{Офіційне видання} \par}
  \vfill
  {\large 2026 \copyright\ ДП «Інформаційні судові системи» \par}
\end{titlepage}

\newpage
\thispagestyle{empty}
\begin{center}
  {\Large\color{nato}\bfseries ЛИСТ ЗАТВЕРДЖЕННЯ ТА КОНТРОЛЮ ВЕРСІЙ\par}
\end{center}
\vspace{20pt}

\begin{longtable}{@{}p{4.0cm}p{10.2cm}@{}}
\toprule
\textbf{Реквізит документа} & \textbf{Значення та відомості} \\
\midrule
\textbf{Назва документа} & Регламент проєктного планування при розробці та супроводі ПЗ ДП «ІСС» \\
\textbf{Об'єкт регулювання} & Процеси управління та розробки програмного забезпечення в ДП «ІСС» \\
\textbf{Версія документа} & v1.0.0 (Офіційна редакція) \\
\textbf{Дата затвердження} & 25.06.2026 \\
\textbf{Статус документа} & Затверджено та введено в дію \\
\textbf{Відповідність стандартам} & ISO 21502:2020, ISO 9001:2015, ISO/IEC/IEEE 12207:2017, ISO/IEC/IEEE 42010:2022, ISO/IEC/IEEE 29148:2018, Постанова КМУ № 205, TOGAF \\
\textbf{Розробники} & Департамент управління проєктами та системної архітектури \\
\textbf{Погоджено} & Архітектурний комітет, Директор з технологій (CTO) \\
\textbf{Затвердив} & Генеральний директор ДП «ІСС» \\
\bottomrule
\end{longtable}

\vspace{20pt}
\begin{center}
  {\large\color{nato}\bfseries Історія змін\par}
\end{center}
\vspace{10pt}
\begin{longtable}{@{}p{1.5cm}p{2.0cm}p{3.0cm}p{7.0cm}@{}}
\toprule
\textbf{Версія} & \textbf{Дата} & \textbf{Автор} & \textbf{Опис внесених змін} \\
\midrule
0.1.0 & 28.05.2026 & Департамент управління проєктами & Первинний проєкт регламенту планування \\
0.9.0 & 15.06.2026 & Архітектурний комітет & Додано відповідність TOGAF та ISO/IEC/IEEE 42010 \\
1.0.0 & 25.06.2026 & Робоча група & Затверджено офіційну версію для впровадження \\
\bottomrule
\end{longtable}

\newpage
\begin{center}
  {\Huge\color{nato}\bfseries Зміст\par}
\end{center}
\vspace{40pt}
\tableofcontents

\newpage

\section*{Анотація}
\addcontentsline{toc}{section}{Анотація}
Цей Регламент визначає єдині методологічні засади, життєвий цикл розробки, правила взаємодії учасників проєктних команд, норми ведення технічної документації та контролю якості для Державного підприємства «Інформаційні судові системи» (ДП «ІСС»). Метою регламенту є структурування внутрішніх процесів розробки та супроводу програмного забезпечення для підвищення прозорості, мінімізації операційних ризиків та успішного проходження міжнародних технічних і фінансових аудитів.

\section*{Методологічна основа та відповідність аудиту}
\addcontentsline{toc}{section}{Методологічна основа та відповідність аудиту}
Регламент розроблено на перетині гнучких практик управління (Agile/Scrum) та класичних міжнародних стандартів системної архітектури і управління проєктами. Відповідність вимогам міжнародного аудиту та державному законодавству забезпечується інтеграцією наступних стандартів та нормативно-правових актів:

\begin{enumerate}
    \item \textbf{TOGAF (The Open Group Architecture Framework):} забезпечує структурування архітектури підприємства за чотирма доменами: бізнес-архітектура, архітектура даних, архітектура додатків та технологічна архітектура. Усі проєктні ініціативи (Epic) проходять попередній аналіз на відповідність цільовій архітектурній моделі ІСС.
    \item \textbf{ISO/IEC/IEEE 42010:2022 (Systems and software engineering --- Architecture description):} стандартизує вимоги до опису архітектури систем, виділяючи окремі архітектурні точки зору (views) та забезпечуючи формальну документацію зв'язків між системними компонентами та бізнес-вимогами.
    \item \textbf{ISO/IEC/IEEE 12207:2017 (Systems and software engineering --- Software life cycle processes):} визначає життєвий цикл програмного забезпечення від формування вимог до виведення з експлуатації. Регламентований у цьому документі workflow у системі JIRA повністю відображає процеси реалізації, інтеграції та верифікації відповідно до цього стандарту.
    \item \textbf{ISO/IEC/IEEE 29148:2018 (Systems and software engineering --- Requirements engineering):} встановлює єдині вимоги до процесів інженерії вимог на всіх етапах життєвого циклу систем, структурує розробку специфікацій та Use Cases.
    \item \textbf{ISO 21502:2020 (Project, programme and portfolio management --- Guidance on project management):} структурує процеси управління ресурсами, ризиками, змінами та якістю результатів у межах спринтів та релізів.
    \item \textbf{ISO 9001:2015 (Quality management systems):} вимагає наявності чітких метрик якості, простежуваності вимог та регулярного аналізу відхилень (наприклад, порівняння запланованого та фактичного часу виконання).
    \item \textbf{Постанова Кабінету Міністрів України від 21 лютого 2025 р. № 205:} визначає обов'язкові вимоги до створення, модернізації, модифікації, розвитку, адміністрування та забезпечення функціонування засобів інформатизації, зокрема щодо обов'язкового використання затверджених національних та міжнародних стандартів при фінансуванні проєктів з державного бюджету.
\end{enumerate}

У межах аудиту регламент гарантує \textbf{розподіл обов'язків (Separation of Duties, SoD)}:
\begin{itemize}
    \item \emph{Product Owner} відповідає за бізнес-вимоги та пріоритети;
    \item \emph{Business Analyst} формалізує вимоги та бізнес-правила у вигляді Use Cases;
    \item \emph{Development Team} здійснює технічну реалізацію без права самостійної зміни вимог;
    \item \emph{QA/Tester} виконує незалежне тестування та підтверджує готовність до релізу;
    \item \emph{Project Manager} контролює дотримання регламенту та управляє ризиками;
    \item \emph{JIRA Administrator} забезпечує технічну неможливість обходу життєвого циклу завдань.
\end{itemize}

\newpage

\section{Загальна інформація (Purpose \& Scope)}

\subsection{Мета документа}
Цей документ встановлює обов'язковий до виконання регламент розробки нового програмного забезпечення (ПЗ) та супроводу існуючих рішень у межах діяльності ДП «ІСС». Регламент визначає ролі учасників, життєвий цикл завдань у системі трекінгу JIRA, правила часового оцінювання, порядок реліз-менеджменту, правила ведення технічної документації в Confluence/Git та механізми безперервного вдосконалення процесів.

\subsection{Принципи та підходи до розробки}
Діяльність проєктних команд ДП «ІСС» базується на наступних фундаментальних принципах:
\begin{itemize}
    \item \textbf{Agile-орієнтація:} розробка здійснюється ітеративно-інкрементним шляхом. Це дозволяє оперативно реагувати на зміни нормативно-правового поля, вимог замовників та технологічних стандартів без втрати керованості процесом.
    \item \textbf{Адаптований Scrum-процес:}
    \begin{itemize}
        \item Базова тривалість ітерації (спринту) становить 2 тижні.
        \item Обов'язковими подіями є: Планування спринту (Sprint Planning), Щоденні наради (Daily Stand-up), Огляд результатів спринту (Sprint Review / Demo) та Ретроспектива спринту (Sprint Retrospective).
    \end{itemize}
    \item \textbf{Ітеративність та поетапна доставка цінності:} кожна ітерація має завершуватися створення потенційно працездатного інкременту системи.
    \item \textbf{Прозорість (Transparency):} поточний стан усіх завдань, оцінки трудовитрат та метрики прогресу є відкритими для учасників процесу та стейкхолдерів через системи JIRA, Confluence та аналітичні дашборди Apache Superset.
\end{itemize}

\subsection{Сфера застосування}
Вимоги цього Регламенту є обов'язковими для:
\begin{itemize}
    \item Усіх довгострокових та короткострокових проєктів з розробки, модернізації та супроводу прикладного ПЗ ДП «ІСС».
    \item Усіх членів внутрішніх команд розробки (штатних працівників) та залучених субпідрядних організацій (вендорів).
    \item Обов'язкового обліку трудовитрат у годинах для кожної окремої задачі в контурі JIRA.
\end{itemize}

\subsection{Взаємодія з іншими політиками та документами}
Цей Регламент діє спільно та узгоджено з Політиками технічних вимог, технічного завдання та технічної документації ІСС, що базуються на стандартах ДСТУ 3973-2000, ДСТУ 3974-2000, а також обов'язкових вимогах до засобів інформатизації, затверджених Постановою Кабінету Міністрів України від 21 лютого 2025 р. № 205 \cite{kmu205}. У разі виникнення розбіжностей між цим Регламентом та вимогами внутрішньої інформаційної безпеки (щодо КСЗІ та НД ТЗІ), пріоритет мають вимоги безпеки.

\section{Ролі та зони відповідальності (Roles \& Responsibilities)}

Для забезпечення прозорості та відповідності вимогам зовнішнього аудиту, ролі в проєктах ДП «ІСС» чітко розмежовані. Спільне виконання ролей розробника та тестувальника, або аналітика та аудитора безпеки в межах одного завдання заборонено.

\subsection{Власник Продукту (Product Owner)}
Власник Продукту є представником бізнес-інтересів замовника (наприклад, ДСА України, органів судової влади) і відповідає за максимізацію цінності продукту.
\begin{itemize}
    \item \textbf{Ключові обов'язки:}
    \begin{itemize}
        \item Створення, наповнення та одноосібна пріоритизація Беклогу Продукту (Product Backlog).
        \item Визначення бізнес-цілей на кожен спринт та погодження складу завдань під час планування.
        \item Тісна співпраця з бізнес-аналітиками для деталізації вимог.
        \item Прийняття або відхилення реалізованого функціоналу на Sprint Review відповідно до Критеріїв Прийняття (Acceptance Criteria).
    \end{itemize}
    \item \textbf{Результат діяльності:} актуальний, прозорий та пріоритизований беклог; чітко сформульовані цілі розробки.
\end{itemize}

\subsection{Керівник Проєкту з функціями Scrum-майстра (Project Manager / Scrum Master)}
Роль об'єднує функції класичного управління проєктними обмеженнями (терміни, бюджети, ризики, звіти) та методологічного лідерства (Scrum Master).
\begin{itemize}
    \item \textbf{Ключові обов'язки:}
    \begin{itemize}
        \item Організація та фасилітація всіх SCRUM-подій, забезпечення дотримання дисципліни ітерацій.
        \item Усунення перешкод (impediments) на шляху команди розробки.
        \item Координація зовнішніх залежностей між суміжними командами та системами.
        \item Управління ресурсами проєкту, бюджетування та комунікація зі стейкхолдерами.
        \item Забезпечення щоденного логування часу командою.
    \end{itemize}
    \item \textbf{Результат діяльності:} ритмічна робота команди без затримок та перешкод; сформовані релізні плани та дотримання бюджетних лімітів.
\end{itemize}

\subsection{Команда Розробки (Development Team)}
Команда є крос-функціональною та містить розробників (Frontend/Backend/Fullstack), тестувальників (QA Manual/Automation), дизайнерів (UI/UX) та технічних письменників.
\begin{itemize}
    \item \textbf{Ключові обов'язки:}
    \begin{itemize}
        \item Декомпозиція та технічна оцінка (у годинах) завдань, включених у спринт.
        \item Безпосередня розробка, тестування, інтеграція та документування функціоналу.
        \item Підтримка актуального стану своїх задач у JIRA протягом дня.
        \item Участь у демонстраціях результатів та ретроспективах.
    \end{itemize}
    \item \textbf{Результат діяльності:} працездатний інкремент продукту в кінці кожного спринту, що повністю відповідає критеріям Definition of Done (DoD).
\end{itemize}

\subsection{Бізнес-аналітик (Business Analyst)}
Бізнес-аналітик залучається з пулу аналітиків ДП «ІСС» на вимогу Project Manager для складних інтеграційних чи масштабних функціональних завдань.
\begin{itemize}
    \item \textbf{Ключові обов'язки:}
    \begin{itemize}
        \item Збір, аналіз та формалізація вимог від Власника Продукту та стейкхолдерів.
        \item Створення специфікацій використання (Use Cases), опис бізнес-правил та побудова BPMN-діаграм процесів.
        \item Трансформація вимог у зрозумілі для команди розробки завдання (JIRA Stories/Tasks).
    \end{itemize}
    \item \textbf{Результат діяльності:} задокументовані Use Cases, BPMN-діаграми та детальні описи бізнес-правил.
\end{itemize}

\subsection{Адміністратор JIRA (JIRA Administrator)}
Окрема технічна роль, що забезпечує стабільність та незмінність життєвого циклу розробки.
\begin{itemize}
    \item \textbf{Ключові обов'язки:}
    \begin{itemize}
        \item Конфігурування проєктів, типів задач, полів та схем переходів (workflow).
        \item Керування правами доступу, забезпечення неможливості ручного переведення задач в обхід обов'язкових статусів (наприклад, перехід з Development відразу в Done без стадії Testing).
        \item Налаштування інтеграцій JIRA з Confluence, репозиторіями коду та Apache Superset.
    \end{itemize}
\end{itemize}

\section{Опис робочих процесів (Workflow Description)}

\subsection{Ієрархія задач і типи робіт у JIRA}
Для забезпечення простежуваності (traceability) від стратегічних цілей до конкретного рядка коду використовується наступна ієрархія сутностей у JIRA:
\begin{enumerate}
    \item \textbf{Theme (Тема):} глобальний стратегічний напрям (наприклад, «Законодавчі зміни 2026», «Міграція на мікросервісну архітектуру»). Реалізується через додавання обов'язкових тегів до Епіків.
    \item \textbf{Epic (Епік):} значний обсяг функціоналу, який не може бути виконаний за один спринт (наприклад, «Модуль виконавчих документів»). Епік має чіткі межі та критерії готовності.
    \item \textbf{Story / Task (Користувацька історія / Задача):} 
    \begin{itemize}
        \item \emph{Story} описує вимогу з точки зору кінцевого користувача та приносить бізнес-цінність.
        \item \emph{Task} описує технічні, інфраструктурні чи дослідницькі роботи (наприклад, «Налаштування індексу бази даних»).
    \end{itemize}
    \item \textbf{Sub-task (Підзадача):} технічна декомпозиція Story/Task на конкретні кроки для окремих виконавців (розробка API, верстка інтефейсу, написання автотестів).
    \item \textbf{Bug (Дефект):} опис невідповідності фактичної поведінки системи вимогам технічного завдання.
    \item \textbf{Release Deployment (Розгортання релізу):} спеціальний системний тип задачі для керування процесом збирання, тестування на Staging та розгортання на Production.
\end{enumerate}

\subsection{Створення й узгодження задач}
Будь-яка задача у JIRA починає свій життєвий цикл з опису. Ініціатор задачі (BA, PO або PM) зобов'язаний заповнити:
\begin{itemize}
    \item Стислий та зрозумілий заголовок.
    \item Детальний опис (для Stories --- за шаблоном: \emph{«Як [Роль] я хочу [Дія] для того, щоб [Цінність]»}).
    \item Критерії Прийняття (Acceptance Criteria), за якими тестувальник перевірятиме результат.
    \item Посилання на пов'язані Use Cases у Confluence або BPMN-діаграми.
\end{itemize}
Задача вважається узгодженою та переводиться в статус \texttt{Ready for Development} тільки після підтвердження Власником Продукту бізнес-цінності та Керівником Проєкту --- повноти опису.

\subsection{Пріоритизація задач}
Пріоритети задач у JIRA розподіляються наступним чином:
\begin{itemize}
    \item \textbf{Highest (Блокуючий):} критичні дефекти на промисловому середовищі (Production), які зупиняють роботу ключових сервісів, або завдання, невиконання яких загрожує безпеці чи порушує вимоги законодавства з визначеним дедлайном. Розробка інших планових задач припиняється для усунення дефекту.
    \item \textbf{High (Високий):} важливі планові завдання або дефекти, що мають обхідні шляхи (workarounds), але суттєво впливають на користувачів.
    \item \textbf{Medium (Середній):} стандартний пріоритет для більшості планових задач та дефектів інтерфейсу.
    \item \textbf{Low (Низький):} другорядні покращення, рефакторинг або косметичні дефекти.
\end{itemize}

\subsection{Життєвий цикл задач (Workflow)}
Для всіх Stories, Tasks та Bugs встановлюється єдиний workflow переходів статусів (див. табл. \ref{tab:workflow}).

\begin{longtable}{@{}p{3.5cm}p{10.7cm}@{}}
\caption{Опис статусів життєвого циклу задач у JIRA} \label{tab:workflow} \\
\toprule
\textbf{Статус у JIRA} & \textbf{Опис процесу та критерії переходу} \\
\midrule
\endfirsthead
\toprule
\textbf{Статус у JIRA} & \textbf{Опис процесу та критерії переходу} \\
\midrule
\endhead
\bottomrule
\endfoot
\bottomrule
\endlastfoot
\texttt{To Do} & Задача створена, перебуває в беклозі. Очікує первинного аналізу та пріоритизації. \\
\texttt{To Analysis} & Затверджено потребу в реалізації. Задача передається аналітику для збору вимог та деталізації. \\
\texttt{Analysis} & Ведеться активне проектування, опис Use Cases, створення макетів та схем. \\
\texttt{To Development} & Критерії готовності до розробки виконані: вимоги задокументовані, погоджені PO, технічна оцінка внесена. \\
\texttt{Development} & Задача в роботі у розробника/дизайнера. Створюється програмний код. \\
\texttt{To Testing} & Код написаний, перевірений (Code Review), розгорнутий на тестовому середовищі. Очікує черги на тестування. \\
\texttt{Testing} & QA-інженер виконує верифікацію на відповідність Критеріям Прийняття. \\
\texttt{For Release testing} & Тестування успішно пройдено. Задача очікує включення в реліз та збірку. \\
\texttt{Done} & Реліз успішно задеплоєно на Production, моніторинг завершено, зауваження відсутні. \\
\end{longtable}

\textbf{Правила циклу доопрацювання (Rework Loop):}
\begin{itemize}
    \item Якщо під час статусу \texttt{Testing} виявлено баги, QA-інженер створює дочірні підзадачі (\emph{Sub-tasks}) типу \texttt{Bug} із детальним описом кроків відтворення.
    \item Основна задача повертається у статус \texttt{Development} (або \texttt{Fix in progress} залежно від налаштувань проєкту).
    \item Після виправлення дефектів розробником, задача знову переводиться у статус \texttt{Testing} для повторної верифікації. Цикл повторюється до повного усунення критичних дефектів.
\end{itemize}

\section{Спринт-процес (Sprint Process)}

\subsection{Планування спринту (Sprint Planning)}
Проводиться у перший день спринту. Тривалість --- до 2 годин.
\begin{enumerate}
    \item \textbf{Вхідні дані:} Беклог Продукту, орієнтовна швидкість команди (Velocity) за минулі спринти, цілі спринту від Власника Продукту.
    \item \textbf{Процес:}
    \begin{itemize}
        \item Обговорення задач зі статусу \texttt{To Development}.
        \item Оцінювання завдань командою розробки в годинах (\emph{Original Estimate}).
        \item На основі доступної ємності команди (Capacity) формирується Спринт Беклог (Sprint Backlog).
        \item Формулюється чітка Ціль Спринту (Sprint Goal).
    \end{itemize}
    \item \textbf{Результат:} зафіксований обсяг задач у спринті, підтверджений зобов'язанням (commitment) команди.
\end{enumerate}

\subsection{Виконання спринту (Sprint Execution)}
\begin{itemize}
    \item \textbf{Daily Stand-up:} щоденна синхронізація тривалістю до 15 хвилин. Кожен виконавець звітує: \emph{що зроблено вчора, що планується сьогодні, які є проблеми/блокери}. PM фіксує блокери та вживає заходів для їх усунення.
    \item \textbf{Логування часу (Worklog):} кожен виконавець зобов'язаний щоденно списувати фактично витрачений час на виконання задач у JIRA.
\end{itemize}

\subsection{Огляд результатів спринту (Sprint Review / Demo)}
Проводиться в останній день спринту. Команда демонструє Власнику Продукту та зацікавленим стейкхолдерам готовий інкрементований функціонал безпосередньо на тестовому або демонстраційному стенді. Презентація слайдів замість демонстрації працюючого коду не допускається. Власник Продукту приймає завдання, які відповідають критеріям прийняття. Неприйняті задачі переносяться в наступний спринт або повертаються в беклог.

\subsection{Ретроспектива спринту (Sprint Retrospective)}
Проводиться після Sprint Review. Мета --- аналіз внутрішніх процесів, інструментів та взаємодії. Результатом є перелік конкретних дій з покращення (Action Items) із визначеними власниками та термінами виконання для наступного спринту.

\section{Реліз-менеджмент (Release Management)}

Реліз-менеджмент ДП «ІСС» спрямований на забезпечення стабільності та безпеки інформаційних систем судової влади при розгортанні змін.

\subsection{Формування переліку задач для релізу}
Керівник Проєкту спільно з Власником Продукту формують релізний пул. 
\begin{itemize}
    \item \textbf{За замовчуванням} до релізу включаються всі задачі, які на момент завершення спринту перебувають у статусі \texttt{For Release testing}, та виправлені дефекти (\texttt{Bugs}) у статусі \texttt{Done}.
    \item Завдання, що мають незакриті підзадачі або блокуючі залежності, виключаються з релізу та переносяться на наступний цикл.
    \item Усі програмні компоненти релізу проходять обов'язкову перевірку на відповідність вимогам локалізації та ліцензійної чистоти відповідно до положень Постанови КМУ № 205 \cite{kmu205}.
\end{itemize}

\subsection{Задача типу Release Deployment}
Для кожного випуску створюється окрема задача типу \texttt{Release Deployment}, яка містить:
\begin{itemize}
    \item Версію релізу відповідно до семантичного версіонування (наприклад, \texttt{v1.4.0}).
    \item План розгортання (Deployment Plan) та перелік відповідальних інженерів.
    \item Перелік пов'язаних задач JIRA, що входять до цієї версії.
    \item Посилання на підписаний QA-інженером звіт про тестування (Test Report) та Release Notes.
\end{itemize}

\subsection{Процес розгортання та моніторингу}
\begin{enumerate}
    \item \textbf{Тестування на Staging:} збірка розгортається на інтеграційному середовищі (Staging) для проведення регресійного та навантажувального тестування.
    \item \textbf{Production Deployment:} після отримання погодження від CTO та PO, здійснюється розгортання на промисловому середовищі (Production) відповідно до затвердженого вікна обслуговування (maintenance window).
    \item \textbf{Моніторинг:} після деплою встановлюється період підвищеної уваги (від 1 до 5 робочих днів). У разі виявлення критичних помилок терміново формується Hotfix.
    \item \textbf{Завершення:} після успішного моніторингу задача \texttt{Release Deployment} переводиться в статус \texttt{Done}, що автоматично переводить усі пов'язані задачі з \texttt{For Release testing} у статус \texttt{Done}.
\end{enumerate}

\section{Управління завданнями та звітність (Reporting)}

\subsection{Відстеження прогресу й трудовитрат}
Для цілей фінансового та технічного аудиту впроваджено жорсткі правила обліку часу:
\begin{itemize}
    \item \textbf{Original Estimate (Первинна оцінка):} вноситься командою під час планування спринту. \textbf{Забороняється змінювати} це поле після старту спринту. Це забезпечує збереження базової лінії (baseline) для оцінки точності планування.
    \item \textbf{Worklog (Журнал робіт):} фактично витрачений час логується виконавцями щодня з точністю до 15 хвилин.
    \item \textbf{Remaining Estimate (Залишок оцінки):} розраховується автоматично або коригується вручну розробником, якщо під час виконання виявлено приховану складність.
\end{itemize}

\subsection{Збирання та візуалізація даних в Apache Superset}
Дані з JIRA (Worklogs, статуси, виконавці, дати створення/закриття) щоденно вивантажуються через ETL-процеси у базу даних аналітики. В Apache Superset налаштовані автоматичні дашборди для керівництва ДП «ІСС» та зовнішніх аудиторів. Звіти будуються за проєктами, епіками, виконавцями та ролями.

\section{Правила оформлення та ведення документації}

Усі проєкти ДП «ІСС» повинні мати актуальну технічну документацію. Відсутність оновлень документації прирівнюється до невиконання вимог задачі.

\subsection{Документація вимог (Use Cases / Specifications)}
Процеси збору, аналізу, формалізації та верифікації вимог до програмного забезпечення здійснюються відповідно до положень стандарту ISO/IEC/IEEE 29148:2018 \cite{iso29148}. Відповідно до вимог Порядку локалізації програмних продуктів для виконання Національної програми інформатизації, затвердженого Постановою Кабінету Міністрів України від 21 лютого 2025 р. № 205 \cite{kmu205}, усі прикладні програмні продукти, інтерфейси користувача та супровідна технічна документація для ЄСІКС повинні виконуватись державною мовою. Для кожної бізнес-функції створюється Use Case у Confluence за затвердженим шаблоном:
\begin{enumerate}
    \item Назва та унікальний ідентифікатор Use Case.
    \item Актори (користувачі, зовнішні системи).
    \item Передумови (Pre-conditions) та Післяумови (Post-conditions).
    \item Основний сценарій взаємодії (кроки користувача та відповіді системи).
    \item Альтернативні та негативні сценарії (обробка помилок, відсутність зв'язку).
    \item BPMN-діаграма для візуалізації складних процесів.
\end{enumerate}
Кожне завдання в JIRA повинно містити пряме посилання на відповідний Use Case.

\subsection{Release Notes}
Створюється для кожного релізу. Містить перелік нових функцій, виправлених багів, опис міграції баз даних та інструкції для адміністраторів щодо розгортання.

\subsection{Технічна документація (Technical Architecture)}
Описує архітектуру системи відповідно до стандартів TOGAF та ISO/IEC/IEEE 42010. Вона включає:
\begin{itemize}
    \item Схеми компонентів та розгортання (Deployment Diagrams).
    \item Специфікацію API методів (запити, відповіді, коди помилок) у форматі OpenAPI/Swagger.
    \item Схеми структур даних (ER-діаграми).
\end{itemize}
Внесення змін до коду без відображення в технічній документації заборонено.

\section{Безперервне вдосконалення процесу (Continuous Improvement)}

\subsection{Регулярні ревізії процесу}
Не рідше ніж один раз на квартал Департамент управління проєктами разом із Архітектурним комітетом проводить ревізію діючих процесів. |Аналізується статистика точності оцінок, швидкість усунення дефектів та зворотний зв'язок від команд.

\subsection{Процедура внесення змін до Регламенту}
Будь-який член команди розробки має право ініціювати зміни до цього Регламенту через пропозицію на Retrospective. Пропозиція розглядається PMO та CTO. У разі затвердження, зміни вносяться до цього документа, оновлюється номер версії, а оновлений регламент публікується в Confluence з обов'язковим Teams-оповіщенням для всіх команд.

\section{Впровадження регламенту та контроль (Implementation \& Controls)}

Контроль за дотриманням цього Регламенту покладається на Керівників Проєктів та внутрішню службу аудиту якості ДП «ІСС». Основні контрольні точки аудиту (Audit Checkpoints) включають:
\begin{enumerate}
    \item Перевірку відповідності архітектурних описів стандарту ISO/IEC/IEEE 42010:2022 перед затвердженням ТЗ.
    \item Перевірку наявності Use Cases та Критеріїв Прийняття у кожній задачі перед переведенням у статус \texttt{To Development}.
    \item Контроль відповідності розгортання релізів вимогам безпеки та наявності підписаних Test Reports.
    \item Перевірку повноти заповнення Worklog у JIRA на щотижневій основі.
\end{enumerate}

\section{Пропоновані аналітичні панелі в Apache Superset}

Для моніторингу ефективності процесів та надання звітів аудиторам в Apache Superset розгортаються наступні візуальні панелі:

\begin{enumerate}
    \item \textbf{Burndown / Burnup Chart (Діаграма згоряння/накопичення):} відображає прогрес виконання годин у межах спринту. Дозволяє виявити ризик зриву цілей спринту на ранніх етапах.
    \item \textbf{Sprint Velocity / Capacity Chart (Швидкість та ємність команди):} показує динаміку виконання годин від спринту до спринту, допомагаючи точніше планувати майбутнє завантаження.
    \item \textbf{Planned vs Actual per Sprint (Заплановані vs Фактичні години):} порівнює сумарний \emph{Original Estimate} із залогованим часом (\emph{Worklog}) після завершення спринту. Використовується для оцінки точності планування.
    \item \textbf{Worklog \& Original Estimate Accuracy (Точність оцінки задач):} звіт по задачах із найбільшим відхиленням фактичних витрат від початкової оцінки для подальшого аналізу на ретроспективі.
    \item \textbf{Task Breakdown by Status / Epic (Розподіл задач за статусами/епіками):} показує поточну концентрацію трудовитрат та виявляє «вузькі місця» (наприклад, накопичення задач у статусі \texttt{To Testing}).
    \item \textbf{Bug Trend \& Bug Severity (Динаміка дефектів):} аналізує кількість відкритих/закритих багів, їх критичність та середній час усунення дефекту (Lead Time for Bugs).
    \item \textbf{Release Status Dashboard (Стан релізів):} відображає готовність наступного релізу, показуючи частку задач у статусі \texttt{For Release testing} та кількість відомих дефектів.
    \item \textbf{Work Distribution by Roles (Розподіл робіт за ролями):} показує баланс завантаженості між аналітиками, розробниками, дизайнерами та тестувальниками.
    \item \textbf{Lead Time / Cycle Time Analysis (Аналіз часу виконання):} вимірює середній час проходження задачі від створення до статусу \texttt{Done} (Lead Time) та від старту розробки до готовності (Cycle Time).
    \item \textbf{Theme Dashboard (Дашборд стратегічних тем):} агрегує витрати часу та прогрес виконання у розрізі глобальних стратегічних ініціатив підприємства.
\end{enumerate}

\newpage

\begin{thebibliography}{99}

\bibitem{togaf} The Open Group.
\newblock \emph{The TOGAF Standard, 10th Edition}.
\newblock The Open Group, 2022.

\bibitem{iso42010} ISO/IEC/IEEE 42010:2022.
\newblock \emph{Systems and software engineering --- Architecture description}.
\newblock International Organization for Standardization, 2022.

\bibitem{iso12207} ISO/IEC/IEEE 12207:2017.
\newblock \emph{Systems and software engineering --- Software life cycle processes}.
\newblock IEEE, 2017.

\bibitem{iso21502} ISO 21502:2020.
\newblock \emph{Project, programme and portfolio management --- Guidance on project management}.
\newblock International Organization for Standardization, 2020.

\bibitem{iso9001} ISO 9001:2015.
\newblock \emph{Quality management systems --- Requirements}.
\newblock International Organization for Standardization, 2015.

\bibitem{dstu3973} ДСТУ 3973-2000.
\newblock \emph{Система розроблення та постачання продукції на виробництво. Правила виконання науково-дослідних робіт. Загальні положення}.
\newblock Держстандарт України, 2000.

\bibitem{dstu3974} ДСТУ 3974-2000.
\newblock \emph{Система розроблення та постачання продукції на виробництво. Правила виконання науково-конструкторських робіт. Загальні положення}.
\newblock Держстандарт України, 2000.

\bibitem{scrumguide} Ken Schwaber and Jeff Sutherland.
\newblock \emph{The Scrum Guide: The Definitive Guide to Scrum: The Rules of the Play}.
\newblock Scrum.Org, 2020.

\bibitem{pmbok} Project Management Institute.
\newblock \emph{A Guide to the Project Management Body of Knowledge (PMBOK Guide), Seventh Edition}.
\newblock Project Management Institute, 2021.

\bibitem{iso15489} ISO 15489-1:2016.
\newblock \emph{Information and documentation --- Records management --- Part 1: Concepts and principles}.
\newblock International Organization for Standardization, 2016.

\bibitem{iso29148} ISO/IEC/IEEE 29148:2018.
\newblock \emph{Systems and software engineering --- Requirements engineering}.
\newblock IEEE, 2018.

\bibitem{kmu205} Постанова Кабінету Міністрів України від 21 лютого 2025 р. № 205.
\newblock \emph{«Деякі питання створення, адміністрування та забезпечення функціонування засобу інформатизації»}.
\newblock Кабінет Міністрів України, 2025.

\end{thebibliography}

\end{document}
